%
% File main.tex
%
% NLP course project-proposal template, Warsaw University of Technology.
% Prepared by Anna Wróblewska, adapted from the ACL 2015 conference
% formatting template (acl2015.sty), originally by the ACL 2015
% organizers and earlier ACL/EACL format authors (see acl2015.pdf
% for the full credits chain).

\documentclass[11pt]{article}
\usepackage{acl2015}
\usepackage{times}
\usepackage{url}
\usepackage{latexsym}
\usepackage{xcolor}

\setlength\titlebox{7.2cm}

% You can expand the title box if you need extra space
% to show all the authors. Please do not make the title box
% smaller than 5cm (the original size); we will check this
% in the camera-ready version and ask you to change it back.


\title{Polish vs. English: measuring the LLM performance gap on matched NLP tasks\\Project Proposal for NLP Course, Winter 2026/27}
% https://www.overleaf.com/project/63440c16442dba647b52a1ae
\author{Piotr Czarnecki \\
  Album number ABCDEF \\
  {\tt 01180766@pw.edu.pl} \\\And
  Emil Łasocha \\
  Album number 341168 \\
  {\tt 01202252@pw.edu.pl} \\ \AND
  Patryk Sukiennik \\
  Album number ABCDEF\\
  {\tt 01180679@pw.edu.pl} \\ \And
  Michał Szewczak \\
  Album number 327320 \\
  {\tt 01180750@pw.edu.pl} \\ \AND
  supervisor: Michał Gromadzki\\
  Warsaw University of Technology \\
  {\tt XXXXXXXXXXX@pw.edu.pl}}

\date{}

\begin{document}
\maketitle

{\small\em Template: Anna Wróblewska (Warsaw University of Technology), adapted from the ACL 2015 conference formatting template.}
\vspace{4pt}

\begin{abstract}
  Large Language Models are evaluated mostly in English, but users also work in different languages such as Polish.
  The proposal of the project consists of building a small and reproducible Polish--English benchmark, evaluating several Large Language Models under identical prompts, and quantifying how much of the observed performance gap is due to model size, prompting and translation quality.
\end{abstract}
\section{Introduction}
Benchmark scores reported for English do not transfer directly to Polish.
A drop in accuracy on a Polish version of a task may reflect a real weakness of the model in Polish, but also artefacts of how the Polish version was produced (machine translation errors, translationese), of how the model is prompted, or of how Polish text is tokenized.
Without separating these factors, a ``Polish gap'' is hard to interpret.
\section{Requirements}
Detailed description (\textbf{max. 4 meaningful pages}) in English, including the following information:
\begin{enumerate}
\item scientific goal of the project (description of the problem to be solved, research questions and hypotheses);
\item significance of the project (state of the art, justification for tackling a specific scientific problem, justification for the pioneering nature of the project, the impact of the project results on the development of the research field and scientific discipline);
\item concept and work plan (general work plan - with a timeline, specific research goals, results of preliminary research, risk analysis);
\item approach \& research methodology (underlying scientific methodology, methods, techniques and research tools, methods of results analysis, equipment and devices to be used in research);
\item \textit{well-justified project literature (a reference list for publications included in the project description, with full bibliographic data);}
\begin{itemize} 
\item BoolQ \cite{clark2019boolqexploringsurprisingdifficulty} evaluates whether a model can answer natural yes/no questions using a supporting passage. It uses approximately 16.000 naturally occurring questions paired with Wikipedia paragraphs. The model receives a question and its passage and predicts either “yes” or “no.” The best reported result is 80.43\% accuracy, achieved by BERT-large with MultiNLI pretraining, fine-tuned on BoolQ. This dataset will be used as one of the benchmarks in our project.

\item GSM8K \cite{DBLP:journals/corr/abs-2110-14168} introduces curated dataset of 8.5 thousand high-quality grade school math problems with natural-language solutions. It then proposes to employ training verifiers to improve LMM's multi-step mathematical reasoning. The generator was fine-tuned for 2 epochs. 100 solutions per training problem were then sampled from it. Each solution was labeled as correct or incorrect, based only on the final answer. Then a separate verifier was trained for 1 epoch to rank solutions. Verifiers help achieve better performance than the same LLM with 30x model size and they scale well with additional data. Additionally, token-level verifiers outperformed solution-level verifiers.


\item MGSM \cite{shi2022languagemodelsmultilingualchainofthought} evaluates the multilingual mathematical reasoning abilities of Large Language Models. It consists of 250 grade--school math problems from the GSM8K dataset, manually translated into ten diverse languages. The model receives a mathematical question in one of the target languages and generates a numerical answer, either directly or by providing intermediate reasoning steps using chain-of-thought prompting. The best reported result is 55.0\% accuracy, achieved by PaLM-540B using English chain-of-thought reasoning after translating the questions into English.
\item Okapi \cite{lai2023okapiinstructiontunedlargelanguage} introduces a framework for instruction-tuning large language models in multiple languages using reinforcement learning from human feedback (RLHF). It covers 26 languages and provides approximately 158.000 instruction examples translated from English. The models are evaluated on their ability to follow instructions and perform knowledge-based and reasoning tasks across different languages. The best reported average result is 49.6\% accuracy on HellaSwag, achieved by the LLaMA-7B model fine-tuned using RLHF.
\item Self-Translate \cite{etxaniz2023multilinguallanguagemodelsthink} investigates whether multilingual language models perform better on reasoning tasks when processing inputs in English rather than in their original language. It introduces a method called self-translate, in which the model first translates the input into English using four-shot prompting and then solves the task using the translated text. The approach is evaluated on five multilingual benchmarks: XStoryCloze, XCOPA, XNLI, PAWS-X, and MGSM, covering common-sense reasoning, natural language inference, paraphrase identification, and mathematical reasoning across multiple languages. Experiments are conducted using XGLM models ranging from 564M to 7.5B parameters and LLaMA models ranging from 7B to 30B parameters. The best reported average result is 46.6\% accuracy, achieved by LLaMA-30B using self-translate, compared to 43.1\% with direct inference. The results show that translating inputs into English consistently improves performance across the evaluated models.


\item XNLI \cite{DBLP:journals/corr/abs-1809-05053} introduces a cross-lingual natural language inference (NLI) benchmark for evaluating cross-lingual sentence representations. It extends the MultiNLI corpora to 15 languages, including low-resource ones. Consisting of 7500 human-annotated development and test examples in NLI three-way classification format, it makes up a total of 112,500 annotated pairs.

\item PLLuM \cite{kocon2025pllum} presents a family of open-source polish foundation models, ranging from 8B to 70B parameters. A 140 billion token polish pretraining corpus was constructed from various texts. Special care was taken to extensively document legal compliance under polish and EU's framework. This model was later used in mObywatel as a virtual assistant.


\item Artetxe et al. \cite{DBLP:journals/corr/abs-2004-04721} shows that human and machine translations can introduce subtle artifacts into data, which cross-lingual models are sensitive to. Fine-tuning on back-translated (round-trip translated) text improves performance on translated test improves accuracy in the translate-test setting by about 4.6 points on average. Training on machine-translated text also improved results, by about 2 points in zero-shot setting. Improvements only apply to translated texts. Based on these results, XNLI state of the art was improved by about 4.3 points in translate-test and 2.8 points in zero-shot.



\end{itemize}
\item \textit{data statement (new for 2026/27) -- source of the data, licence, whether it contains personal data, and what you do about it. Six lines are enough;}
\begin{table*}[t]
\centering
\small
\setlength{\tabcolsep}{4pt}
\renewcommand{\arraystretch}{1.25}
\begin{tabular}{p{1.2cm} p{2.5cm} p{2.5cm} p{1.9cm} p{1.7cm} p{2cm} p{2.2cm}}
\hline
\textbf{Dataset} &
\textbf{Size} &
\textbf{Label definition} &
\textbf{Class balance} &
\textbf{Licence} &
\textbf{How obtained} &
\textbf{Status} \\
\hline

MMLU
& 14,042 test questions
& Multiple-choice answers (4 options)
& To verify
& To verify
& Public benchmark
& READY + TRANSLATION TO COLLECT \\

GSM8K
& 8,792 test questions
& Numerical answers
& Not applicable
& To verify
& Public benchmark
& READY + TRANSLATION TO COLLECT \\

ARC-Challenge
& 2,590 test questions
& Multiple-choice answers
& To verify
& To verify
& Public benchmark
& READY + TRANSLATION TO COLLECT \\

BoolQ
& 9,427 train records and 3,270 validation records
& Question, Paragraph, Answer
& Majority of True label, 62.31\%
& cc-by-sa-3.0
& Public dataset
& READY + TRANSLATION TO COLLECT \\

\hline
\end{tabular}
\caption{Overview of the datasets.}
\label{tab:datasets}
\end{table*}
As shown in Table~\ref{tab:datasets}, we consider four popular English benchmarks for our study. The original English datasets are publicly available -- READY, while Polish translations are TO COLLECT. Table~\ref{tab:datasets} also lists the licences under which the benchmarks are distributed. The BoolQ dataset does not contain any personal information, as its answers are based on Wikipedia paragraphs \cite{clark2019boolqexploringsurprisingdifficulty}. A subset of each benchmark will be translated using machine translation and reviewed.

\item \textbf{contribution table} -- each team member, what they did, workload in hours. \textbf{Required in this and every later phase report; a report submitted without it is not accepted for grading, not just marked down};
\item \textbf{declaration of LLM use} -- which tools, for what. \textbf{A report without this is not accepted for grading} (new for 2026/27).
\end{enumerate}

\textbf{Reminder: place this and every later submission in a folder named \texttt{surname\_surname\_surname\_surname} (all team members). The order set here must not change for the rest of the semester -- a mismatch costs 5 points.}

% include your bib file like this:
%\bibliographystyle{acl}
%\bibliography{acl2015}
\bibliographystyle{acl}
\bibliography{bibliography}
\end{document}
